% !TEX root = ../main.tex
% =====================================================================
% 中文摘要
%
% 以下文字与附件 3 中文摘要页完全一致，逐字照录（含附件 3 自带的格式说明
% 段落）。示例中“摘要：”之后为说明文字，实际使用时请替换为论文摘要正文
% （规范要求不少于 400 字）。
%
% 环境参数：#1 = 摘要页题目（黑体三号，居中，段前 1 行、段后 0.5 行）
%           #2 = 关键词（词间用中文分号“；”，末位不加标点）
% =====================================================================

\begin{cnabstract}{中文论文题目}{关键词1；关键词2；关键词3；关键词4；关键词5}

{\heiti\zihao{4}摘要}：摘要是论文的内容不加注释和评论的简短陈述。摘要应具有独立性和自含性，即不阅读论文的全文，就能获得必要的信息。摘要中有数据、有结论，是一篇完整的短文，可以独立使用。摘要应说明研究工作目的、实验方法、结果和最终结论等，重点是结果和结论。中文摘要一般不少于400字。如遇特殊需要，字数可略多。摘要中一般不用图、表、化学结构式、非公知公用的符号和术语。

摘要论文题目黑体三号，分1或2行居中打印，段前间距为1行，段后间距为0.5行；论文题目下“专业班级、学生姓名、指导教师、职称”，居中，宋体小四号，1.5倍行距。

“摘要”二字用黑体四号，1.5倍行距，段前0.5行，首行缩进2字符。“摘要”二字后打印摘要内容，内容用宋体小四号，1.5倍行距，首行缩进2字符。

关键词是为了文献标引工作从论文中选取出来的以表示全文主题内容信息款目的单词或术语。每篇论文选取3-5个词作为关键词，以显著的字符另起一行，排在摘要的左下方。尽量用《汉语主题词表》等词表提供的规范词。

关键词段前间距为1行，左顶格。“关键词”三字用黑体四号，“关键词”后为内容，用宋体小四号，每两个关键词之间用中文分号“；”间隔，最后一个关键词后不加标点符号。

从原创性声明开始设置页眉，原创性声明、中文摘要、英文摘要、目录页的页眉分两栏，左侧以小五号宋体分别键入“原创性声明”、“摘要”、“ABSTRACT”、“目录”，右侧统一键入“江苏大学本科毕业设计（论文）”。

从正文开始，以小五号宋体键入“江苏大学本科毕业设计（论文）”，靠右侧显示。原创性声明、中文摘要、英文摘要、目录页的页码使用罗马数字（Ⅰ、Ⅱ、Ⅲ）编排；正文页码从绪论开始至附录（正文-致谢-参考文献-附录）使用阿拉伯数字编排，小五号Times New Roman居中显示。

\end{cnabstract}
